% Options for packages loaded elsewhere
\PassOptionsToPackage{unicode}{hyperref}
\PassOptionsToPackage{hyphens}{url}
\documentclass[
]{article}
\usepackage{xcolor}
\usepackage{amsmath,amssymb}
\setcounter{secnumdepth}{-\maxdimen} % remove section numbering
\usepackage{iftex}
\ifPDFTeX
  \usepackage[T1]{fontenc}
  \usepackage[utf8]{inputenc}
  \usepackage{textcomp} % provide euro and other symbols
\else % if luatex or xetex
  \usepackage{unicode-math} % this also loads fontspec
  \defaultfontfeatures{Scale=MatchLowercase}
  \defaultfontfeatures[\rmfamily]{Ligatures=TeX,Scale=1}
\fi
\usepackage{lmodern}
\ifPDFTeX\else
  % xetex/luatex font selection
\fi
% Use upquote if available, for straight quotes in verbatim environments
\IfFileExists{upquote.sty}{\usepackage{upquote}}{}
\IfFileExists{microtype.sty}{% use microtype if available
  \usepackage[]{microtype}
  \UseMicrotypeSet[protrusion]{basicmath} % disable protrusion for tt fonts
}{}
\makeatletter
\@ifundefined{KOMAClassName}{% if non-KOMA class
  \IfFileExists{parskip.sty}{%
    \usepackage{parskip}
  }{% else
    \setlength{\parindent}{0pt}
    \setlength{\parskip}{6pt plus 2pt minus 1pt}}
}{% if KOMA class
  \KOMAoptions{parskip=half}}
\makeatother
\usepackage{color}
\usepackage{fancyvrb}
\newcommand{\VerbBar}{|}
\newcommand{\VERB}{\Verb[commandchars=\\\{\}]}
\DefineVerbatimEnvironment{Highlighting}{Verbatim}{commandchars=\\\{\}}
% Add ',fontsize=\small' for more characters per line
\newenvironment{Shaded}{}{}
\newcommand{\AlertTok}[1]{\textcolor[rgb]{1.00,0.00,0.00}{\textbf{#1}}}
\newcommand{\AnnotationTok}[1]{\textcolor[rgb]{0.38,0.63,0.69}{\textbf{\textit{#1}}}}
\newcommand{\AttributeTok}[1]{\textcolor[rgb]{0.49,0.56,0.16}{#1}}
\newcommand{\BaseNTok}[1]{\textcolor[rgb]{0.25,0.63,0.44}{#1}}
\newcommand{\BuiltInTok}[1]{\textcolor[rgb]{0.00,0.50,0.00}{#1}}
\newcommand{\CharTok}[1]{\textcolor[rgb]{0.25,0.44,0.63}{#1}}
\newcommand{\CommentTok}[1]{\textcolor[rgb]{0.38,0.63,0.69}{\textit{#1}}}
\newcommand{\CommentVarTok}[1]{\textcolor[rgb]{0.38,0.63,0.69}{\textbf{\textit{#1}}}}
\newcommand{\ConstantTok}[1]{\textcolor[rgb]{0.53,0.00,0.00}{#1}}
\newcommand{\ControlFlowTok}[1]{\textcolor[rgb]{0.00,0.44,0.13}{\textbf{#1}}}
\newcommand{\DataTypeTok}[1]{\textcolor[rgb]{0.56,0.13,0.00}{#1}}
\newcommand{\DecValTok}[1]{\textcolor[rgb]{0.25,0.63,0.44}{#1}}
\newcommand{\DocumentationTok}[1]{\textcolor[rgb]{0.73,0.13,0.13}{\textit{#1}}}
\newcommand{\ErrorTok}[1]{\textcolor[rgb]{1.00,0.00,0.00}{\textbf{#1}}}
\newcommand{\ExtensionTok}[1]{#1}
\newcommand{\FloatTok}[1]{\textcolor[rgb]{0.25,0.63,0.44}{#1}}
\newcommand{\FunctionTok}[1]{\textcolor[rgb]{0.02,0.16,0.49}{#1}}
\newcommand{\ImportTok}[1]{\textcolor[rgb]{0.00,0.50,0.00}{\textbf{#1}}}
\newcommand{\InformationTok}[1]{\textcolor[rgb]{0.38,0.63,0.69}{\textbf{\textit{#1}}}}
\newcommand{\KeywordTok}[1]{\textcolor[rgb]{0.00,0.44,0.13}{\textbf{#1}}}
\newcommand{\NormalTok}[1]{#1}
\newcommand{\OperatorTok}[1]{\textcolor[rgb]{0.40,0.40,0.40}{#1}}
\newcommand{\OtherTok}[1]{\textcolor[rgb]{0.00,0.44,0.13}{#1}}
\newcommand{\PreprocessorTok}[1]{\textcolor[rgb]{0.74,0.48,0.00}{#1}}
\newcommand{\RegionMarkerTok}[1]{#1}
\newcommand{\SpecialCharTok}[1]{\textcolor[rgb]{0.25,0.44,0.63}{#1}}
\newcommand{\SpecialStringTok}[1]{\textcolor[rgb]{0.73,0.40,0.53}{#1}}
\newcommand{\StringTok}[1]{\textcolor[rgb]{0.25,0.44,0.63}{#1}}
\newcommand{\VariableTok}[1]{\textcolor[rgb]{0.10,0.09,0.49}{#1}}
\newcommand{\VerbatimStringTok}[1]{\textcolor[rgb]{0.25,0.44,0.63}{#1}}
\newcommand{\WarningTok}[1]{\textcolor[rgb]{0.38,0.63,0.69}{\textbf{\textit{#1}}}}
\usepackage{longtable,booktabs,array}
\newcounter{none} % for unnumbered tables
\usepackage{calc} % for calculating minipage widths
% Correct order of tables after \paragraph or \subparagraph
\usepackage{etoolbox}
\makeatletter
\patchcmd\longtable{\par}{\if@noskipsec\mbox{}\fi\par}{}{}
\makeatother
% Allow footnotes in longtable head/foot
\IfFileExists{footnotehyper.sty}{\usepackage{footnotehyper}}{\usepackage{footnote}}
\makesavenoteenv{longtable}
\setlength{\emergencystretch}{3em} % prevent overfull lines
\providecommand{\tightlist}{%
  \setlength{\itemsep}{0pt}\setlength{\parskip}{0pt}}
\usepackage{bookmark}
\IfFileExists{xurl.sty}{\usepackage{xurl}}{} % add URL line breaks if available
\urlstyle{same}
% fallback for those not using the hyperref driver hyperxmp:
\makeatletter
\@ifundefined{xmpquote}{\newcommand{\xmpquote}[1]{#1}}{}
\makeatother
\hypersetup{
  hidelinks,
  pdfcreator={LaTeX via pandoc}}

\author{}
\date{}

\begin{document}

\section{Manuscript (draft --- export to PDF /
LaTeX)}\label{manuscript-draft-export-to-pdf-latex}

\textbf{Purpose:} Submission-facing prose only. Scope:
\url{paper-roadmap.md}. Claims: \textbf{\texttt{.cursor/memory.md}} and
\url{claim-evidence-map.md}.

\textbf{Last updated:} 2026-05-29.

\textbf{PDF:}
\texttt{pandoc\ docs/paper/manuscript.md\ -o\ paper-draft.pdf\ -\/-pdf-engine=tectonic}.
IOS: \texttt{./docs/paper/submission/build\_ios\_package.sh}.

\begin{center}\rule{0.5\linewidth}{0.5pt}\end{center}

\subsection{Working title}\label{working-title}

\emph{Point-in-Time Multi-Agent LLM Trading: A Reproducible LangGraph
Stack with Regime-Conditioned Case Study Evaluation and Fair
Backtesting}

\begin{center}\rule{0.5\linewidth}{0.5pt}\end{center}

\subsection{Abstract}\label{abstract}

Large language models are increasingly used in financial decision
pipelines, yet published agent stacks are often hard to reproduce,
rarely enforce temporal data boundaries, and seldom report net returns
after transaction costs. We present \textbf{Hindsight 20/20}, an open
\textbf{LangGraph} stack with structured \textbf{Pydantic} outputs, a
\textbf{simulation end date} for point-in-time data, optional
\textbf{ticker anonymization}, and a fee-aware \textbf{paper backtest}.
The primary contribution is \textbf{reproducible systems
engineering}---one codepath, auditable stage artifacts, and frozen CSV
schedules---not a universe-scale alpha benchmark. As an
\textbf{illustrative case study}, we evaluate the \textbf{full} pipeline
on two liquid large-cap NSE listings---\textbf{RELIANCE.NS} and
\textbf{TCS.NS}---in \textbf{four pre-specified regime windows} (two
tickers × exogenous bear/bull labels; \textbf{one deterministic
trajectory per window}, no seed averaging). In \textbf{both bear
windows}, the agent loses \textbf{less} than buy-and-hold on total
return and shows \textbf{lower path drawdown} (RELIANCE: −7.5\% vs
−17.4\%, MDD 5.1\% vs 21.0\%; TCS: −4.8\% vs −20.4\%, MDD 6.8\% vs
23.9\%). In \textbf{bull windows}, outcomes \textbf{diverge}: RELIANCE
lags buy-and-hold (+5.6\% vs +15.1\%) while TCS nearly tracks it
(+19.6\% vs +21.9\%). We interpret heterogeneous bull outcomes as an
\textbf{architectural boundary} of defensive multi-stage coordination
rather than a performance claim. Frozen CSV artifacts and scripts
reproduce tables and figures. \textbf{Code and artifacts:}
\url{https://github.com/RMartires/hindsight}.

\emph{(IOS submission:
\href{submission/}{\texttt{docs/paper/submission/}} --- run
\texttt{build\_ios\_package.sh}.)}

\begin{center}\rule{0.5\linewidth}{0.5pt}\end{center}

\subsection{1. Introduction}\label{introduction}

Researchers are experimenting with large language models (LLMs) inside
trading workflows amid a rapid expansion of LLM-driven agent systems
(Bădică et al., 2025). Yet evaluations rarely ask how the \textbf{same
implementation} behaves when the \textbf{underlying trend} differs,
rarely enforce point-in-time data access at the adapter layer, and
rarely report \textbf{net} returns after stated transaction costs.

\textbf{What this paper is.} \textbf{Hindsight 20/20} extends the
open-source \textbf{TradingAgents} multi-agent LangGraph framework (Xiao
et al., 2024) with (i) schema-constrained stage outputs, (ii) a
simulation end date for point-in-time vendor queries, (iii) optional
ticker anonymization, and (iv) a scriptable, fee-aware paper backtest
tied to \textbf{frozen CSV artifacts}. We document \textbf{how} to build
and audit such a stack and provide an \textbf{honest, scoped}
regime-split case study on two names.

\textbf{What this paper is not.} We do \textbf{not} claim
state-of-the-art risk-adjusted returns, regime \emph{prediction}, or
statistical generalization from a broad equity universe. We do
\textbf{not} report multi-seed variance (each regime window is
\textbf{N=1}). Frontier empirical papers (e.g., Jeon and Lee, 2026;
Miyazaki et al., 2026) average tens of random seeds across index
universes; our evaluation scope is deliberately narrower and is labeled
as such (§4.1, Table 3).

\textbf{Case-study evaluation.} We run the \textbf{full} pipeline
(\texttt{PAPER\_ABLATION=full}) on \textbf{RELIANCE.NS} and
\textbf{TCS.NS}. For each ticker we pre-specify one \textbf{bear} and
one \textbf{bull} calendar window in \textbf{\texttt{regime-timeranges}}
before any agent run; labels follow exogenous underlying buy-and-hold
thresholds (§4.2). Each of \textbf{four regime windows} starts from
\textbf{\$100,000} fresh paper cash; we compare agent \textbf{net}
return to \textbf{buy-and-hold} on the same window.

\textbf{Code.} Implementation, frozen evaluation CSVs, and figure
scripts: \url{https://github.com/RMartires/hindsight}.

\textbf{Structure.} §2 related work; §3 system; §4 case study
evaluation; §5 conclusion; \textbf{Contributions} at the end.

\begin{center}\rule{0.5\linewidth}{0.5pt}\end{center}

\subsection{2. Related Work}\label{related-work}

\emph{Bibliography: \url{bib-seed.bib}.}

\textbf{LLM agents and classic MAS.} Multi-agent frameworks orchestrate
specialist roles for trading and finance (Xiao et al., 2024; Yu et al.,
2023; Yang et al., 2023). Bădică et al.~(2025) situate this wave as a
shift from explicitly programmed BDI-style agents to emergent LLM
behaviors---raising new demands for verifiable coordination and audit
trails. While \textbf{TradingAgents} (Xiao et al., 2024) provides the
foundational analyst → debate → trader graph that we extend, it does not
ship adapter-level temporal clamps, schema-persisted stage artifacts, or
a fee-aware replication harness. \textbf{AlphaAgents} (Zhao et al.,
2025) validates institution-style role separation (Fundamental,
Sentiment, Valuation) with structured debate; \textbf{QuantAgents} (Li
et al., 2025) emphasizes manager-led \textbf{meetings} and
simulated-trading feedback---both comparators for our Bull/Bear debate
loop, but neither foregrounds point-in-time vendor routing or frozen
schedule CSVs as first-class outputs.

\textbf{Methodological standards, leakage, and overfitting.} Nguyen and
Pham (2026) survey LLM financial multi-agent systems and propose
\textbf{Consolidated Minimum Standards}, documenting pervasive
evaluation failures---look-ahead bias, transaction-cost neglect, and
regime-shift blindness---that can reverse the sign of reported returns.
They introduce the \textbf{Coordination Breakeven Spread (CBS)} for
asking whether multi-agent coordination pays for its turnover cost; we
report a simplified fee-and-turnover breakeven in §4.5 and defer full
latency CBS to future work. Bailey et al.~(2014) show that undisclosed
configuration search makes strong in-sample backtests statistically
cheap to manufacture; Hindsight responds with \textbf{pre-specified}
windows in \texttt{regime-timeranges}, frozen artifacts, and explicit
N=1 scope rather than implied generalization.

\textbf{Point-in-time policy and anonymization.} Early LLM trading
papers often relied on prompt-level ``do not peek'' instructions. Nguyen
and Pham (2026) treat \textbf{contamination control} and
\textbf{net-of-cost returns} as minimum standards; Hindsight implements
the former via a \textbf{simulation end date} at vendor routing
(\textbf{M4}) and the latter via a \textbf{Zerodha delivery} paper
ledger. Jeon and Lee (2026) (\textbf{BlindTrade}) provide the strongest
anonymization-first defense: tickers like ``AAPL'' can trigger memorized
pretraining associations rather than fresh reasoning; they anonymize
S\&P 500 constituents daily and report Sharpe 1.40 ± 0.22 across
\textbf{20 seeds}. Hindsight's optional
\textbf{\texttt{STOCK\_\#\#\#\#}} aliases (§3.5) pursue the same goal
within a LangGraph replication stack; we do not ablate anonymization in
§4.

\textbf{Fine-grained tasks and coordination topology.} Miyazaki et
al.~(2026) show that \textbf{fine-grained} analyst tasks (e.g., explicit
ratio and indicator computations) outperform coarse ``analyze
fundamentals'' prompts under leakage-controlled TOPIX-style evaluation
with \textbf{50 trials}. Our market and fundamentals analysts are
\textbf{tool-bound} to concrete fetches and Tier-1 indicators
(§3.2)---aligning architecture with that finding without claiming
comparable universe-scale statistics. Dewansh (2026) frames
\textbf{explainable stability} and \textbf{rationale persistence} as
audit criteria for orchestrated trading agents; our Pydantic stage
outputs and wide schedule CSVs are designed to support that traceability
(outlook → stance → \texttt{final\_signal}), even though we do not
formally score policy convergence.

\textbf{Benchmarks, simulators, and live evaluation.} Qian et al.~(2026)
introduce the \textbf{Agent Market Arena (AMA)}, a lifelong multi-market
benchmark showing that \textbf{agent architecture} often matters more
than LLM backbone choice---a motivation for holding the model fixed
(\texttt{qwen/qwen3.5-flash-02-23}) and varying regime context in §4.
Papadakis et al.~(2025) (\textbf{StockSim}) define order-level
simulation with latency, slippage, and limit-order-book microstructure;
our paper ledger is candlestick-level with stated fees and \textbf{no
LOB claim}, treating StockSim-class realism as aspirational future work.

\textbf{Empirical scope vs frontier benchmarks.} Jeon and Lee (2026), Li
et al.~(2025), and Miyazaki et al.~(2026) optimize for
\textbf{statistical stress tests} across index universes, seeds, or
simulated meetings. Hindsight optimizes for \textbf{auditability and
temporal integrity} on a two-ticker, four-window case study (Table
3)---complementary, not directly comparable on seeds or universe
breadth.

\textbf{Positioning summary.} Where \textbf{TradingAgents} supplies the
graph skeleton and \textbf{BlindTrade} / Nguyen and Pham (2026)
articulate \emph{what to measure}, Hindsight 20/20 implements \emph{how
to reproduce it}: simulation end dates, structured artifacts, optional
anonymization, net returns, and frozen CSV schedules that regenerate
Tables 1--4 without rerunning LLM inference.

\begin{center}\rule{0.5\linewidth}{0.5pt}\end{center}

\subsection{3. Method / System}\label{method-system}

Pointers: \textbf{\texttt{.cursor/memory.md}}.

\subsubsection{3.1 Architecture}\label{architecture}

\textbf{\texttt{TradingAgentsGraph}}: analysts → \textbf{Bull/Bear
Researcher} debate (\emph{internal roles}, not §4 market regimes) →
trader → risk → final decision. Core routing follows
\textbf{TradingAgents} (Xiao et al., 2024); our modifications focus on
temporal data policy, schema-constrained outputs, and the
backtest/evaluation pipeline below. Same graph for
\textbf{\texttt{run\_backtest\_mvp}}.

\subsubsection{3.2 Analyst task
decomposition}\label{analyst-task-decomposition}

Each analyst role is bound to \textbf{specific tools} and emits a
structured \textbf{\texttt{AnalystReport}} (outlook + narrative). This
supports traceable rationale from raw data to downstream debate---not
coarse role labels alone.

{\def\LTcaptype{none} % do not increment counter
\begin{longtable}[]{@{}
  >{\raggedright\arraybackslash}p{(\linewidth - 6\tabcolsep) * \real{0.1273}}
  >{\raggedright\arraybackslash}p{(\linewidth - 6\tabcolsep) * \real{0.1455}}
  >{\raggedright\arraybackslash}p{(\linewidth - 6\tabcolsep) * \real{0.4364}}
  >{\raggedright\arraybackslash}p{(\linewidth - 6\tabcolsep) * \real{0.2909}}@{}}
\toprule\noalign{}
\begin{minipage}[b]{\linewidth}\raggedright
Stage
\end{minipage} & \begin{minipage}[b]{\linewidth}\raggedright
Agent
\end{minipage} & \begin{minipage}[b]{\linewidth}\raggedright
Tools (representative)
\end{minipage} & \begin{minipage}[b]{\linewidth}\raggedright
Concrete tasks
\end{minipage} \\
\midrule\noalign{}
\endhead
\bottomrule\noalign{}
\endlastfoot
Market & Market analyst & \texttt{get\_stock\_data},
\texttt{get\_indicators} & Fetch OHLCV; select up to \textbf{8
complementary Tier-1 indicators} (trend: 50/200 SMA, 10 EMA, MACD;
momentum: RSI; volatility: Bollinger, ATR; volume: VWMA, MFI) with
non-redundant justification \\
Fundamentals & Fundamentals analyst & \texttt{get\_fundamentals},
\texttt{get\_balance\_sheet}, \texttt{get\_cashflow},
\texttt{get\_income\_statement} & Pull statements and issuer profile;
analyze revenue, margins, balance-sheet trends \\
News & News analyst & \texttt{get\_news}, \texttt{get\_global\_news} &
Company-targeted and macro news over a rolling window \\
Sentiment & Social-media analyst & \texttt{get\_news} (company/social
queries) & Sentiment and social discussion synthesis \\
Debate & Bull / Bear researchers & --- (read analyst reports) & Argue
buy vs sell thesis from shared evidence \\
Execution & Trader, Risk manager & --- & Propose and gate
\textbf{BUY/HOLD/SELL} \\
\end{longtable}
}

Structured fields (\texttt{market\_outlook},
\texttt{fundamentals\_outlook}, \texttt{bull\_implied\_stance},
\texttt{final\_signal}, etc.) are persisted in schedule CSVs for audit
(\textbf{M3}), supporting \textbf{explainable stability} in the sense of
Dewansh (2026): a coherent trace from evidence to executable action.

\subsubsection{3.3 Tools and structured
outputs}\label{tools-and-structured-outputs}

Analyst tool calls respect the simulation end date; stage outputs
conform to \textbf{Pydantic} schemas
(\texttt{tradingagents/schemas/outputs.py}, \textbf{M3}).

\subsubsection{3.4 Point-in-time data
policy}\label{point-in-time-data-policy}

A \textbf{simulation end date} clamps vendor queries so agents cannot
access data after the as-of trade date (\textbf{M4}).

\subsubsection{3.5 Ticker anonymization}\label{ticker-anonymization}

LLMs trained on public text encode firm-specific narratives tied to
ticker symbols. In a point-in-time backtest, that \textbf{pretraining
leakage} can bias decisions.

Hindsight optionally \textbf{anonymizes the issuer} for all LLM-facing
prompts and tool outputs (\textbf{M5}):

\begin{enumerate}
\def\labelenumi{\arabic{enumi}.}
\tightlist
\item
  \textbf{Deterministic ticker aliases} ---
  \texttt{TickerMapper.for\_real\_ticker()} maps each real symbol to a
  stable synthetic id (\texttt{STOCK\_\#\#\#\#}, SHA-256--derived).
\item
  \textbf{Round-trip vendor access} --- tool wrappers call
  \texttt{resolve\_ticker\_for\_vendor()} before API calls, then
  \texttt{scrub\_ticker\_text()} on responses.
\item
  \textbf{News scrubbing} --- \texttt{scrub\_news\_text()} replaces
  multi-word proper nouns with \texttt{PROPER\_NOUN} placeholders
  (heuristic, news-only).
\end{enumerate}

Frozen E1 uses \textbf{\texttt{enable\_anonymization=True}}. We do
\textbf{not} ablate anonymization in §4.

\subsubsection{3.6 Backtest metrics}\label{backtest-metrics}

The paper ledger (\texttt{tradingagents/backtest/}) records daily agent
equity \textbf{after fees}. Primary §4 metrics emphasize \textbf{total
return}, \textbf{maximum drawdown (MDD)}, and \textbf{fee drag};
risk-adjusted ratios are secondary on short windows (67--111 trading
days).

{\def\LTcaptype{none} % do not increment counter
\begin{longtable}[]{@{}
  >{\raggedright\arraybackslash}p{(\linewidth - 2\tabcolsep) * \real{0.4000}}
  >{\raggedright\arraybackslash}p{(\linewidth - 2\tabcolsep) * \real{0.6000}}@{}}
\toprule\noalign{}
\begin{minipage}[b]{\linewidth}\raggedright
Metric
\end{minipage} & \begin{minipage}[b]{\linewidth}\raggedright
Definition
\end{minipage} \\
\midrule\noalign{}
\endhead
\bottomrule\noalign{}
\endlastfoot
\textbf{Total return (net)} & ((V\_T - V\_0) / V\_0) on agent equity
after fees \\
\textbf{Gross return} & Net return plus fees as a fraction of (V\_0) \\
\textbf{Fee drag} & Gross return minus net return \\
\textbf{Alpha vs B\&H} & Agent net return minus buy-and-hold return on
the same adjusted closes \\
\textbf{Max drawdown (MDD)} & Peak-to-trough on agent equity;
\textbf{B\&H MDD} on (V\_0 \times \text{close}/\text{close}\_0) \\
\textbf{Excess Sharpe ratio} & Mean daily excess return (over risk-free)
/ stdev, × (\sqrt{252})\footnote{Annualization uses 252 trading days (US
  convention). NSE calendars are ≈240--250 days; we retain 252 for
  comparability. \textbf{Risk-free rate:} India short-rate proxy
  \textbf{6.5\% p.a.}, converted to a daily rate (r\_f/252), applied
  uniformly across windows (post-hoc on frozen equity series).
  Qualitative conclusions are unchanged under a ±1\% p.a. sensitivity
  band.} \\
\textbf{Sortino / Calmar} & Implemented; reported in Table 2
(secondary) \\
\end{longtable}
}

\subsubsection{3.7 Pipeline configuration}\label{pipeline-configuration}

§4 uses \textbf{\texttt{PAPER\_ABLATION=full}}. Presets
\textbf{\texttt{a1}--\texttt{a3}} exist for internal use only
(\textbf{M2}).

\subsubsection{3.8 Paper backtest and replication
packet}\label{paper-backtest-and-replication-packet}

\textbf{\texttt{scripts/backtest\_mvp.py\ -\/-dates-csv}} writes wide
schedule CSVs (signals, equity, outlook, stance, and running metrics).

\textbf{Replication packet (Frozen E1):} four schedule CSVs under
\textbf{\texttt{results/dates-*-dates.csv}}, matching logs under
\textbf{\texttt{eval\_results/}}, pre-specified windows in
\textbf{\texttt{regime-timeranges}}, shared env
(\texttt{PAPER\_ABLATION=full}, \texttt{initial\_cash=100000},
\texttt{cost\_model=zerodha\_delivery}, \texttt{slippage\_bps=0},
anonymization on, LLM \texttt{qwen/qwen3.5-flash-02-23}). Tables and
figures regenerate via
\textbf{\texttt{scripts/backtest\_regime\_analysis.ipynb}} and
\textbf{\texttt{scripts/generate\_paper\_figures.py}}.

\subsubsection{3.9 Observability}\label{observability}

Langfuse optional; not empirical evidence.

\begin{center}\rule{0.5\linewidth}{0.5pt}\end{center}

\subsection{4. Case Study Evaluation}\label{case-study-evaluation}

\subsubsection{4.1 Evaluation scope and inferential
limits}\label{evaluation-scope-and-inferential-limits}

{\def\LTcaptype{none} % do not increment counter
\begin{longtable}[]{@{}
  >{\raggedright\arraybackslash}p{(\linewidth - 4\tabcolsep) * \real{0.2157}}
  >{\raggedright\arraybackslash}p{(\linewidth - 4\tabcolsep) * \real{0.2353}}
  >{\raggedright\arraybackslash}p{(\linewidth - 4\tabcolsep) * \real{0.5490}}@{}}
\toprule\noalign{}
\begin{minipage}[b]{\linewidth}\raggedright
Dimension
\end{minipage} & \begin{minipage}[b]{\linewidth}\raggedright
This study
\end{minipage} & \begin{minipage}[b]{\linewidth}\raggedright
Typical frontier benchmark
\end{minipage} \\
\midrule\noalign{}
\endhead
\bottomrule\noalign{}
\endlastfoot
\textbf{Universe} & 2 liquid NSE large-caps (RELIANCE, TCS) & Index
constituents (e.g., S\&P 500, TOPIX 100) \\
\textbf{Regime windows} & 4 pre-specified (2 tickers × bear/bull) & Full
sample or rolling panels \\
\textbf{Stochasticity} & \textbf{N=1 trajectory per window} (no seed
variance) & 20--50 independent trials \\
\textbf{Regime labels} & Exogenous from underlying B\&H (±10\%) & Often
endogenous or undisclosed \\
\textbf{Bias control} & Simulation end date + optional anonymization &
Strict anonymization (BlindTrade) \\
\textbf{Costs} & Stated delivery fees (net returns) & Often gross
returns \\
\end{longtable}
}

We treat §4 as an \textbf{illustrative stress test} of one stack under
contrasting underlying trends---not a hypothesis test about expected
alpha. Outcomes \textbf{observed in these four windows} should not be
extrapolated to a broad universe without further study (§5).

\textbf{Pre-specification.} Calendar windows are listed in
\textbf{\texttt{regime-timeranges}} \textbf{before} agent backtests.
Regime labels are assigned \textbf{exogenously} from underlying
buy-and-hold return and are \textbf{independent} of internal Bull/Bear
Researcher debate roles.

\subsubsection{4.2 Regime definition and
protocol}\label{regime-definition-and-protocol}

{\def\LTcaptype{none} % do not increment counter
\begin{longtable}[]{@{}
  >{\raggedright\arraybackslash}p{(\linewidth - 2\tabcolsep) * \real{0.2000}}
  >{\raggedright\arraybackslash}p{(\linewidth - 2\tabcolsep) * \real{0.8000}}@{}}
\toprule\noalign{}
\begin{minipage}[b]{\linewidth}\raggedright
Regime label
\end{minipage} & \begin{minipage}[b]{\linewidth}\raggedright
Rule (underlying B\&H on processed window, per ticker)
\end{minipage} \\
\midrule\noalign{}
\endhead
\bottomrule\noalign{}
\endlastfoot
\textbf{Bear} & Total return \textbf{≤ −10\%} \\
\textbf{Bull} & Total return \textbf{≥ +10\%} \\
\end{longtable}
}

Windows \textbf{differ by ticker} (e.g., TCS bull Jun--Sep 2024;
RELIANCE bull Jan--Jun 2025). Each run uses \textbf{\$100,000} initial
cash.

{\def\LTcaptype{none} % do not increment counter
\begin{longtable}[]{@{}
  >{\raggedright\arraybackslash}p{(\linewidth - 6\tabcolsep) * \real{0.1667}}
  >{\raggedright\arraybackslash}p{(\linewidth - 6\tabcolsep) * \real{0.1667}}
  >{\raggedright\arraybackslash}p{(\linewidth - 6\tabcolsep) * \real{0.3750}}
  >{\raggedright\arraybackslash}p{(\linewidth - 6\tabcolsep) * \real{0.2917}}@{}}
\toprule\noalign{}
\begin{minipage}[b]{\linewidth}\raggedright
Ticker
\end{minipage} & \begin{minipage}[b]{\linewidth}\raggedright
Regime
\end{minipage} & \begin{minipage}[b]{\linewidth}\raggedright
Processed window
\end{minipage} & \begin{minipage}[b]{\linewidth}\raggedright
Schedule CSV
\end{minipage} \\
\midrule\noalign{}
\endhead
\bottomrule\noalign{}
\endlastfoot
RELIANCE.NS & Bear & 2024-08-01 .. 2025-01-03 &
\texttt{results/dates-REL-aug-2024-jan-2025-bear-dates.csv} \\
RELIANCE.NS & Bull & 2025-01-01 .. 2025-06-03 &
\texttt{results/dates-REL-jan-2025-june-2025-bull-dates.csv} \\
TCS.NS & Bear & 2024-12-02 .. 2025-04-03 &
\texttt{results/dates-TCS-dec-2024-april-2025-bear-dates.csv} \\
TCS.NS & Bull & 2024-06-03 .. 2024-09-03 &
\texttt{results/dates-TCS-june-2024-sept-2024-bull-dates.csv} \\
\end{longtable}
}

Driver (example):

\begin{Shaded}
\begin{Highlighting}[]
\ExtensionTok{python}\NormalTok{ scripts/backtest\_mvp.py }\AttributeTok{{-}{-}ticker}\NormalTok{ TCS.NS }\DataTypeTok{\textbackslash{}}
  \AttributeTok{{-}{-}dates{-}csv}\NormalTok{ results/dates{-}TCS{-}dec{-}2024{-}april{-}2025{-}bear{-}dates.csv}
\end{Highlighting}
\end{Shaded}

\textbf{Frozen artifacts:} \url{claim-evidence-map.md} (\textbf{Frozen
E1}).

\subsubsection{4.3 Positioning vs prior agent-trading
systems}\label{positioning-vs-prior-agent-trading-systems}

\textbf{Table 3.} Qualitative comparison (see §4.1 for scope limits).

{\def\LTcaptype{none} % do not increment counter
\begin{longtable}[]{@{}
  >{\raggedright\arraybackslash}p{(\linewidth - 8\tabcolsep) * \real{0.1184}}
  >{\raggedright\arraybackslash}p{(\linewidth - 8\tabcolsep) * \real{0.2237}}
  >{\raggedright\arraybackslash}p{(\linewidth - 8\tabcolsep) * \real{0.1579}}
  >{\raggedright\arraybackslash}p{(\linewidth - 8\tabcolsep) * \real{0.1711}}
  >{\raggedright\arraybackslash}p{(\linewidth - 8\tabcolsep) * \real{0.3289}}@{}}
\toprule\noalign{}
\begin{minipage}[b]{\linewidth}\raggedright
Feature
\end{minipage} & \begin{minipage}[b]{\linewidth}\raggedright
Hindsight 20/20
\end{minipage} & \begin{minipage}[b]{\linewidth}\raggedright
BlindTrade
\end{minipage} & \begin{minipage}[b]{\linewidth}\raggedright
QuantAgents
\end{minipage} & \begin{minipage}[b]{\linewidth}\raggedright
Expert Investment Teams
\end{minipage} \\
\midrule\noalign{}
\endhead
\bottomrule\noalign{}
\endlastfoot
Primary goal & Reproducible PIT stack + audit trail & Anonymized RL/GNN
policy & Simulated agent meetings & Hierarchical analyst teams \\
Architecture & LangGraph (staged) & GNN + RL & Meeting simulation &
Manager--analyst hierarchy \\
Stochasticity reported & N=1 per regime window & 20 seeds & Multi-year +
live tracking & 50 trials \\
Universe & 2 NSE large-caps & S\&P 500 (PIT) & Multi-asset simulation &
TOPIX 100 \\
Bias control & Sim. end date + opt. anonymization & Strict anonymization
& Simulation feedback & Knowledge cutoff controls \\
Costs in backtest & Net (Zerodha delivery) & Varies & Varies & Varies \\
Key §4 metrics & Return, MDD, fee drag & Sharpe (mean ± std) & ARR,
volatility & Median Sharpe \\
\end{longtable}
}

\subsubsection{4.4 Results}\label{results}

\textbf{Table 1.} Agent vs buy-and-hold (frozen CSVs, \$100,000 start).
\textbf{Primary metrics:} net return, MDD, alpha vs B\&H, fees. Returns
and MDD recomputed from daily \texttt{equity} and \texttt{close} columns
in frozen CSVs.

{\def\LTcaptype{none} % do not increment counter
\begin{longtable}[]{@{}
  >{\raggedright\arraybackslash}p{(\linewidth - 22\tabcolsep) * \real{0.0769}}
  >{\raggedright\arraybackslash}p{(\linewidth - 22\tabcolsep) * \real{0.0769}}
  >{\raggedright\arraybackslash}p{(\linewidth - 22\tabcolsep) * \real{0.0577}}
  >{\raggedright\arraybackslash}p{(\linewidth - 22\tabcolsep) * \real{0.0481}}
  >{\raggedright\arraybackslash}p{(\linewidth - 22\tabcolsep) * \real{0.1250}}
  >{\raggedright\arraybackslash}p{(\linewidth - 22\tabcolsep) * \real{0.0673}}
  >{\raggedright\arraybackslash}p{(\linewidth - 22\tabcolsep) * \real{0.0962}}
  >{\raggedright\arraybackslash}p{(\linewidth - 22\tabcolsep) * \real{0.0962}}
  >{\raggedright\arraybackslash}p{(\linewidth - 22\tabcolsep) * \real{0.0865}}
  >{\raggedright\arraybackslash}p{(\linewidth - 22\tabcolsep) * \real{0.1058}}
  >{\raggedright\arraybackslash}p{(\linewidth - 22\tabcolsep) * \real{0.0962}}
  >{\raggedright\arraybackslash}p{(\linewidth - 22\tabcolsep) * \real{0.0673}}@{}}
\toprule\noalign{}
\begin{minipage}[b]{\linewidth}\raggedright
Ticker
\end{minipage} & \begin{minipage}[b]{\linewidth}\raggedright
Regime
\end{minipage} & \begin{minipage}[b]{\linewidth}\raggedright
Days
\end{minipage} & \begin{minipage}[b]{\linewidth}\raggedright
B\&H
\end{minipage} & \begin{minipage}[b]{\linewidth}\raggedright
Agent (net)
\end{minipage} & \begin{minipage}[b]{\linewidth}\raggedright
Gross
\end{minipage} & \begin{minipage}[b]{\linewidth}\raggedright
Fee drag
\end{minipage} & \begin{minipage}[b]{\linewidth}\raggedright
α vs B\&H
\end{minipage} & \begin{minipage}[b]{\linewidth}\raggedright
B\&H MDD
\end{minipage} & \begin{minipage}[b]{\linewidth}\raggedright
Agent MDD
\end{minipage} & \begin{minipage}[b]{\linewidth}\raggedright
Fees (\$)
\end{minipage} & \begin{minipage}[b]{\linewidth}\raggedright
B/H/S
\end{minipage} \\
\midrule\noalign{}
\endhead
\bottomrule\noalign{}
\endlastfoot
RELIANCE & Bear & 111 & −17.43\% & −7.54\% & −6.53\% & 1.01 pp & +9.89
pp & 21.0\% & 5.11\% & 1,010 & 7/28/76 \\
RELIANCE & Bull & 108 & +15.09\% & +5.56\% & +6.19\% & 0.64 pp & −9.53
pp & 11.0\% & 2.80\% & 635 & 9/25/74 \\
TCS & Bear & 88 & −20.42\% & −4.83\% & −2.92\% & 1.91 pp & +15.59 pp &
23.9\% & 6.78\% & 1,910 & 12/24/52 \\
TCS & Bull & 67 & +21.86\% & +19.61\% & +21.80\% & 2.20 pp & −2.25 pp &
5.5\% & 1.65\% & 2,198 & 13/14/40 \\
\end{longtable}
}

\textbf{Table 2.} Secondary risk-adjusted metrics (post-hoc on frozen
daily equity; excess Sharpe uses \textbf{6.5\% p.a.} risk-free
rate\footnote{Annualization uses 252 trading days (US convention). NSE
  calendars are ≈240--250 days; we retain 252 for comparability.
  \textbf{Risk-free rate:} India short-rate proxy \textbf{6.5\% p.a.},
  converted to a daily rate (r\_f/252), applied uniformly across windows
  (post-hoc on frozen equity series). Qualitative conclusions are
  unchanged under a ±1\% p.a. sensitivity band.}).

{\def\LTcaptype{none} % do not increment counter
\begin{longtable}[]{@{}
  >{\raggedright\arraybackslash}p{(\linewidth - 10\tabcolsep) * \real{0.1053}}
  >{\raggedright\arraybackslash}p{(\linewidth - 10\tabcolsep) * \real{0.1053}}
  >{\raggedright\arraybackslash}p{(\linewidth - 10\tabcolsep) * \real{0.2895}}
  >{\raggedright\arraybackslash}p{(\linewidth - 10\tabcolsep) * \real{0.2763}}
  >{\raggedright\arraybackslash}p{(\linewidth - 10\tabcolsep) * \real{0.1184}}
  >{\raggedright\arraybackslash}p{(\linewidth - 10\tabcolsep) * \real{0.1053}}@{}}
\toprule\noalign{}
\begin{minipage}[b]{\linewidth}\raggedright
Ticker
\end{minipage} & \begin{minipage}[b]{\linewidth}\raggedright
Regime
\end{minipage} & \begin{minipage}[b]{\linewidth}\raggedright
Excess Sharpe (agent)
\end{minipage} & \begin{minipage}[b]{\linewidth}\raggedright
Excess Sharpe (B\&H)
\end{minipage} & \begin{minipage}[b]{\linewidth}\raggedright
Sortino
\end{minipage} & \begin{minipage}[b]{\linewidth}\raggedright
Calmar
\end{minipage} \\
\midrule\noalign{}
\endhead
\bottomrule\noalign{}
\endlastfoot
RELIANCE & Bear & −4.27 & −2.51 & −3.10 & −0.96 \\
RELIANCE & Bull & +0.74 & +1.25 & +3.67 & +1.94 \\
TCS & Bear & −2.69 & −3.13 & −2.25 & −0.71 \\
TCS & Bull & +3.86 & +3.27 & +24.22 & +11.88 \\
\end{longtable}
}

Short windows and discrete daily rebalancing make risk ratios
\textbf{noisy}; we interpret Table 2 alongside Table 1, not in
isolation. RELIANCE bull is positive on total return (+5.6\%) with
positive excess Sharpe (+0.74) but lags B\&H on alpha (−9.53 pp).

\textbf{Cross-ticker patterns (observed in these windows).}

\begin{itemize}
\tightlist
\item
  \textbf{Bear (2/2):} agent \textbf{loses less} than B\&H on total
  return and shows \textbf{markedly lower MDD} (e.g., RELIANCE bear:
  5.1\% vs 21.0\%). Both runs are \textbf{SELL-heavy}, consistent with
  reduced long exposure during drawdowns.
\item
  \textbf{Bull (mixed):} TCS agent \textbf{nearly tracks} B\&H (+19.6\%
  vs +21.9\%, α −2.25 pp); RELIANCE agent \textbf{lags} (+5.6\% vs
  +15.1\%, α −9.53 pp) with 74 SELL vs 9 BUY days.
\end{itemize}

\textbf{Table 4.} Selected outlook / stance shares (\% of trading days).
Full pivot in notebook.

{\def\LTcaptype{none} % do not increment counter
\begin{longtable}[]{@{}
  >{\raggedright\arraybackslash}p{(\linewidth - 10\tabcolsep) * \real{0.0964}}
  >{\raggedright\arraybackslash}p{(\linewidth - 10\tabcolsep) * \real{0.0964}}
  >{\raggedright\arraybackslash}p{(\linewidth - 10\tabcolsep) * \real{0.1928}}
  >{\raggedright\arraybackslash}p{(\linewidth - 10\tabcolsep) * \real{0.1928}}
  >{\raggedright\arraybackslash}p{(\linewidth - 10\tabcolsep) * \real{0.2048}}
  >{\raggedright\arraybackslash}p{(\linewidth - 10\tabcolsep) * \real{0.2169}}@{}}
\toprule\noalign{}
\begin{minipage}[b]{\linewidth}\raggedright
Ticker
\end{minipage} & \begin{minipage}[b]{\linewidth}\raggedright
Regime
\end{minipage} & \begin{minipage}[b]{\linewidth}\raggedright
market bullish
\end{minipage} & \begin{minipage}[b]{\linewidth}\raggedright
market bearish
\end{minipage} & \begin{minipage}[b]{\linewidth}\raggedright
bull stance buy
\end{minipage} & \begin{minipage}[b]{\linewidth}\raggedright
bear stance sell
\end{minipage} \\
\midrule\noalign{}
\endhead
\bottomrule\noalign{}
\endlastfoot
RELIANCE & Bear & 10.8 & 31.5 & 100.0 & 88.3 \\
RELIANCE & Bull & 24.1 & 19.4 & 94.4 & 77.8 \\
TCS & Bear & 22.7 & 33.0 & 71.6 & 53.4 \\
TCS & Bull & 46.3 & 14.9 & 64.2 & 44.8 \\
\end{longtable}
}

Internal \textbf{bull researcher} stances can remain buy-leaning while
\textbf{final signals} stay defensive---downstream trader and risk
stages drive executable actions.

\subsubsection{4.5 Transaction-cost breakeven and architectural
boundaries}\label{transaction-cost-breakeven-and-architectural-boundaries}

\textbf{Fee drag vs alpha.} Fee drag is modest (0.6--2.2 pp of initial
capital; Table 1). It does \textbf{not} explain RELIANCE bull
under-participation: the \textbf{−9.53 pp alpha gap} dominates
\textbf{0.64 pp} fee drag. In bear windows, defensive positioning yields
\textbf{positive alpha} (+9.89 pp REL, +15.59 pp TCS) despite similar
fee levels---suggesting \textbf{exposure management}, not costs, drives
regime-split outcomes.

\textbf{Simplified breakeven spread.} Let (K) be non-HOLD signal days
and (F) total fees. A coarse \textbf{fee-and-turnover breakeven} is the
average daily edge required to justify rebalancing: (F / (V\_0
\cdot K)). For RELIANCE bull ((K=83) actionable days, (F=\$635)), this
is ≈\textbf{7.7 bps/day}---far below the \textbf{\textasciitilde88
bps/day} implied alpha shortfall ((9.53\%/108) days). The binding
constraint is \textbf{SELL-heavy coordination} (74 SELL days), not
delivery fees alone.

\textbf{Architectural boundary (hypothesis).} Multi-stage debate →
trader → risk can \textbf{compress drawdown} in bear windows but
\textbf{cap rally participation} in bull windows when final policy stays
defensive despite buy-leaning researcher stances (Table 4). RELIANCE
bull exemplifies this boundary; TCS bull shows the same stack can
\textbf{nearly track} B\&H when signal balance is less skewed (13 BUY,
40 SELL). We frame this as \textbf{regime × issuer interaction under one
architecture}, not a universal win/loss verdict.

We do \textbf{not} estimate latency-driven slippage (full
\textbf{Coordination Breakeven Spread} per Nguyen and Pham, 2026); that
remains future work alongside StockSim-class order-level simulation
(Papadakis et al., 2025).

\subsubsection{4.6 Figures}\label{figures}

\begin{Shaded}
\begin{Highlighting}[]
\VariableTok{MPLCONFIGDIR}\OperatorTok{=}\NormalTok{.mpl\_cache }\ExtensionTok{.venv/bin/python}\NormalTok{ scripts/generate\_paper\_figures.py}
\end{Highlighting}
\end{Shaded}

\begin{itemize}
\tightlist
\item
  \textbf{\texttt{fig\_regime\_equity\_grid}} --- 2×2 agent vs B\&H
  panels\\
\item
  \textbf{\texttt{fig\_regime\_returns\_by\_ticker}} --- total return
  bars (four runs)\\
\item
  \textbf{\texttt{fig\_regime\_signals\_by\_ticker}} --- BUY/HOLD/SELL
  counts\\
\item
  \textbf{\texttt{fig\_regime\_outlook\_bullish}} --- bullish
  analyst-outlook share heatmap
\end{itemize}

\begin{center}\rule{0.5\linewidth}{0.5pt}\end{center}

\subsection{5. Conclusion}\label{conclusion}

We presented \textbf{Hindsight 20/20}, a reproducible point-in-time
LangGraph stack with structured outputs, optional anonymization, and
fee-aware backtesting, plus a \textbf{scoped four-window case study} on
\textbf{RELIANCE.NS} and \textbf{TCS.NS}. In the \textbf{observed bear
windows}, the full pipeline shows \textbf{lower loss and drawdown} than
buy-and-hold; \textbf{bull-window outcomes are heterogeneous},
consistent with an architectural trade-off between defensive
coordination and rally participation.

\textbf{Limitations and future work.}

\begin{itemize}
\tightlist
\item
  \textbf{Inferential scope:} N=1 per regime window; no multi-seed
  variance; two tickers---not an index universe.
\item
  \textbf{Calendars:} Per-ticker windows are not synchronous; no shared
  macro calendar claim.
\item
  \textbf{Model and ablations:} Single LLM
  (\texttt{qwen/qwen3.5-flash-02-23}); anonymization on but not ablated;
  \texttt{a1}--\texttt{a3} presets not evaluated in §4.
\item
  \textbf{Metrics:} Short windows; risk ratios are secondary; no
  latency/slippage CBS; Infosys windows in \texttt{regime-timeranges}
  planned but not run.
\item
  \textbf{Future empirical work (out of scope here):} multi-seed trials,
  broader PIT universes, anonymization ablations, and full CBS with
  measured agent latency.
\end{itemize}

Regimes are \textbf{exogenous labels}, not agent predictions.
\textbf{Not claimed:} SOTA returns; regime forecasting; statistical
generalization beyond the four frozen windows.

\begin{center}\rule{0.5\linewidth}{0.5pt}\end{center}

\subsection{Contributions}\label{contributions}

\begin{enumerate}
\def\labelenumi{\arabic{enumi}.}
\tightlist
\item
  \textbf{Reproducible backtest harness and frozen artifacts} ---
  fee-aware paper ledger, schedule CSVs, replication scripts (S5 /
  M6).\\
\item
  \textbf{Point-in-time data policy and optional anonymization} ---
  simulation end date at vendor routing; deterministic ticker aliases
  (S4 / M4--M5).\\
\item
  \textbf{Structured stage outputs} --- Pydantic-constrained artifacts
  for audit and traceability (S3 / M3).\\
\item
  \textbf{Multi-agent decision graph} with tool-bound analyst
  decomposition (S1 / M1).\\
\item
  \textbf{Honest regime-conditioned case study} --- four pre-specified
  windows, documented limits, architectural-boundary interpretation (S2
  / E1).
\end{enumerate}

\textbf{Not claimed:} SOTA returns; regime prediction; multi-seed
benchmarks; ablation studies in §4.

\begin{center}\rule{0.5\linewidth}{0.5pt}\end{center}

\subsection{Acknowledgements}\label{acknowledgements}

The agent graph implementation builds on \textbf{TradingAgents} (Xiao et
al., 2024; Tauric Research, Apache License 2.0). We thank the authors
for open-sourcing the framework.

\begin{center}\rule{0.5\linewidth}{0.5pt}\end{center}

\subsection{References}\label{references}

Badica, C., Badica, A., Ganzha, M., Ivanović, M., Paprzycki, M.,
Selişteanu, D. and Wrona, Z., 2025. Contemporary Agent Technology:
LLM-Driven Advancements vs Classic Multi-Agent Systems. \emph{arXiv
preprint} arXiv:2509.02515 {[}cs.MA{]}. https://arxiv.org/abs/2509.02515

Bailey, D.H., Borwein, J.M., López de Prado, M. and Zhu, Q.J., 2014.
Pseudo-Mathematics and Financial Charlatanism: The Effects of Backtest
Overfitting on Out-of-Sample Performance. \emph{Notices of the American
Mathematical Society}, 61(5), pp.458--471.
https://doi.org/10.1090/noti1105

Dewansh, S., 2026. Policy Convergence and Explainable Stability in
Orchestrated Multi-Agent LLM Trading Frameworks. \emph{International
Journal of Advanced Engineering and Management Science}, 12(2),
pp.058--064. https://doi.org/10.22161/ijaems.122.8

Jeon, J. and Lee, H., 2026. Can Blindfolded LLMs Still Trade? An
Anonymization-First Framework for Portfolio Optimization. \emph{arXiv
preprint} arXiv:2603.17692 {[}cs.LG{]}. https://arxiv.org/abs/2603.17692

Li, X., Zeng, Y., Xing, X., Xu, J. and Xu, X., 2025. QuantAgents:
Towards Multi-agent Financial System via Simulated Trading. In
\emph{Findings of the Association for Computational Linguistics: EMNLP
2025}, pp.17438--17464, Suzhou, China.
https://doi.org/10.18653/v1/2025.findings-emnlp.945

Miyazaki, K., Kawahara, T., Roberts, S. and Zohren, S., 2026. Toward
Expert Investment Teams: A Multi-Agent LLM System with Fine-Grained
Trading Tasks. \emph{arXiv preprint} arXiv:2602.23330 {[}cs.AI{]}.
https://arxiv.org/abs/2602.23330

Nguyen, P. and Pham, T., 2026. Toward Reliable Evaluation of LLM-Based
Financial Multi-Agent Systems: Taxonomy, Coordination Primacy, and Cost
Awareness. \emph{arXiv preprint} arXiv:2603.27539 {[}cs.MA{]}.
https://arxiv.org/abs/2603.27539

Papadakis, C., Filandrianos, G., Dimitriou, A., Lymperaiou, M., Thomas,
K. and Stamou, G., 2025. StockSim: A Dual-Mode Order-Level Simulator for
Evaluating Multi-Agent LLMs in Financial Markets. \emph{arXiv preprint}
arXiv:2507.09255 {[}cs.CE{]}. https://arxiv.org/abs/2507.09255

Qian, L., Peng, X., Wang, Y., Zhang, V.J., He, H., Li, Z., Smith, H.,
Han, Y., He, Y., Li, H., Cao, Y., Yu, Y., Lopez-Lira, A., Lu, P., Nie,
J.-Y., Xiong, G., Huang, J. and Ananiadou, S., 2026. When Agents Trade:
Live Multi-Market Trading Arena for LLM Agents. In \emph{Proceedings of
the ACM Web Conference 2026}. https://doi.org/10.1145/3774904.3792821

Xiao, Y., Sun, E., Luo, D. and Wang, W., 2024. TradingAgents:
Multi-Agents LLM Financial Trading Framework. \emph{arXiv preprint}
arXiv:2412.20138 {[}q-fin.TR{]}. https://arxiv.org/abs/2412.20138

Yang, H., Liu, X.-Y. and Wang, C.D., 2023. FinGPT: Open-Source Financial
Large Language Models. \emph{arXiv preprint} arXiv:2306.06031
{[}q-fin.ST{]}. https://arxiv.org/abs/2306.06031

Yu, Y., Li, H., Chen, Z., Jiang, Y., Li, Y., Zhang, D., Liu, R., Suchow,
J.W. and Khashanah, K., 2023. FinMem: A Performance-Enhanced LLM Trading
Agent with Layered Memory and Character Design. \emph{arXiv preprint}
arXiv:2311.13743 {[}q-fin.CP{]}. https://arxiv.org/abs/2311.13743

Zhao, T., Lyu, J., Jones, S., Garber, H., Pasquali, S. and Mehta, D.,
2025. AlphaAgents: Large Language Model based Multi-Agents for Equity
Portfolio Constructions. \emph{arXiv preprint} arXiv:2508.11152
{[}q-fin.ST{]}. https://arxiv.org/abs/2508.11152

\begin{center}\rule{0.5\linewidth}{0.5pt}\end{center}

\subsection{Changelog}\label{changelog}

{\def\LTcaptype{none} % do not increment counter
\begin{longtable}[]{@{}
  >{\raggedright\arraybackslash}p{(\linewidth - 2\tabcolsep) * \real{0.4286}}
  >{\raggedright\arraybackslash}p{(\linewidth - 2\tabcolsep) * \real{0.5714}}@{}}
\toprule\noalign{}
\begin{minipage}[b]{\linewidth}\raggedright
Date
\end{minipage} & \begin{minipage}[b]{\linewidth}\raggedright
Change
\end{minipage} \\
\midrule\noalign{}
\endhead
\bottomrule\noalign{}
\endlastfoot
2026-05-29 & Expanded references and Related Work contrasts (Nguyen and
Pham, Jeon and Lee, Bailey et al., Miyazaki et al., Li et al., Zhao et
al., Papadakis et al., Qian et al., Dewansh, Bădică et al.). \\
2026-05-29 & Multi-ticker Frozen E1: TCS + completed REL bull through
2025-06-03. \\
\end{longtable}
}

\end{document}
